\subsection{策梅洛定理}

引理3. 设E是一个集合，设S是 \(\mathfrak{P}(\mathbf{E})\) 的一个子集，并且令 \(p\) 为从S到E中的映射，使得 \(p(\mathbf{X}) \notin \mathbf{X}\) （对于所有 \(\mathbf{X} \in \mathfrak{S}\) ）。那么存在E的一个子集M和M上的一个良序 \(\Gamma\) ，使得，如果 \(x \leqslant y\) 表示关系 \(y \in \Gamma\langle x \rangle\) 并且 \(\mathrm{S}_x\) 表示段 {]}←, \(x]\) ,

\begin{enumerate}
\def\labelenumi{(\arabic{enumi})}
\item
  对于所有 \(x \in \mathbf{M}\) 我们有 \(\mathbf{S}_x \in \mathfrak{S}\) 并且 \(p(\mathbf{S}_x) = x\) ;
\item
  \(\mathbf{M} \notin \mathfrak{S}\) .
\end{enumerate}

令 \(\mathfrak{M}\) 为由满足以下条件的子集 \(\mathbf{G}\) （\(\mathbf{E} \times \mathbf{E}\) 的子集）所组成的集合：

\begin{enumerate}
\def\labelenumi{(\alph{enumi})}
\item
  \(\mathbf{G}\) 是 \(\mathbf{pr}_1\mathbf{G} = \mathbf{U}\) 上的一个良序的图;
\item
  如果 \(x \leqslant y\) 表示关系 \((x, y) \in \mathbf{G}\) （在 \(\mathbf{U}\) 上），那么对于每个 \(x \in \mathbf{U}\) ，段 \(S_x\) 满足 \(S_x \in \mathfrak{S}\) 并且 \(p(S_x) = x\) 。我们将证明，如果 \(G\) 与 \(G'\) 是 \(M\) 的两个元素，并且 \(U, U'\) 表示它们的第一投影，则这两个集合 \(U, U'\) 中的一个包含在另一个之中，并且如果，例如， \(U \subset U'\) ，那么 \(G = G' \cap (U \times U)\) （换句话说， \(U\) 上的序关系是由 \(U'\) 上的序关系诱导的）且 \(U\) 是 \(U'\) 的一个段。
\end{enumerate}

考虑集合 V，其元素 \(x \in U \cap U'\) 满足：以 x 为端点的段在 U 和 \(U'\) 中相同，并且该段上由 U 和 \(U'\) 上的序关系所诱导的序关系也相同。显然，V 在 U 和 \(U'\) 中都是段，并且在V上诱导的序关系相同；因此，只要我们证明要么V = U，要么 \(V = U'\) ，断言即得证。让我们用反证法来论证，并假设 \(V \neq U\) 并且 \(V \neq U'\) 。设 x 是 U - V 在 U 中的最小元素，并设 \(x'\) 是 \(U' - V\) 在 \(U'\) 中的最小元素；我们有 \(V = S_x\) （在 U 中），且 \(V = S_x'\) （在 \(U'\) 中）。但根据假设， \(V \in S\) 并且

\[
x = p \left(\mathrm{S} _ {x}\right), \quad x ^ {\prime} = p \left(\mathrm{S} _ {x ^ {\prime}}\right),
\]

使得 \(x = x'\) 。因此根据定义 \(x \in V\) ，这是荒谬的。因此我们可以将第1号中的命题3应用于第一投影的集合 \(U = pr_1G\) （其中 \(G \in \mathfrak{M}\) ），从而得到一个良序集

\[
\mathrm{M} = \bigcup_ {\mathrm{G} \in \mathfrak {M}} \operatorname{pr} _ {1} \mathrm{G}.
\]

容易看出，M上的序关系的图属于M。如果我们有 \(M \in S\) ，则令 \(a = p(M)\) ，我们应当有 \(a \notin M\) . 因此我们可以将元素 a 作为最大元素添加到 M 中，并且集合 \(M' = M \cup \{a\}\) 将是良序的。由于 \(M = S_a\) （在 \(M'\) 中），我们应当有 \(S_a \in S\) 并且 \(p(S_a) = a\) ； \(M'\) 上的序关系的图因此将属于M，这是荒谬的。

注意，如果 \(\phi \notin \mathfrak{S}\) （特别是如果 \(\mathfrak{S}\) 为空），则由引理3断言其存在的集合M是空集；这由引理3的条件1得出。

\textbf{定理1（策梅洛）。每个集合E都可以被良序化。}

令 \(\mathfrak{S} = \mathfrak{P}(\mathrm{E}) - \{\mathrm{E}\}\) 为 \(\mathbf{E}\) 的除 \(\mathbf{E}\) 自身之外的所有子集所成的集合。对于每个 \(\mathbf{X} \in \mathfrak{S}\) 令 \(p(\mathbf{X}) = \tau_x\) ( \(x \in \mathbf{E} - \mathbf{X}\) ）；由于关系 \(\mathbf{X} \in \mathfrak{S}\) 蕴含 \((\exists x)(x \in \mathbf{E} - \mathbf{X})\) ，我们有 \(p(\mathbf{X}) \in \mathbf{E} - \mathbf{X}\) （第一章，§4，第1号），因此 \(p(\mathbf{X}) \notin \mathbf{X}\) 。因此我们可以应用引理3，从而在 \(\mathbf{E}\) 的一个子集 \(\mathbf{M}\) （使得 \(\mathbf{M} \notin \mathfrak{S}\) ）上存在一个良序；但 \(\mathbf{E}\) 的唯一不属于 \(\mathfrak{S}\) 的子集是 \(\mathbf{E}\) 本身，定理得证。

